\begin{table}
\begin{tcolorbox}[tabularx={X|p{2cm}|X|},enhanced,fonttitle=\bfseries\large,fontupper=\normalsize\sffamily,
colback=yellow!10!white,colframe=red!50!black,colbacktitle=blue!30!white,
coltitle=black,title=Miscellaneous Orbiter Commands (Part 2)]
\hline
Command & Arguments & Purpose \\
\hline
\hline
\verb'-change_values' & fname-in fname-out map-in map-out & 
Change the values of the file in fname-in according to a mapping. 
The output is written to fname-out.
The mapping is defined by two corresponding sequences. 
The input values are in map-in, 
the corresponding output values are in map-out. \\
\hline  
\verb'-store_as_csv_file' & fname $m$ $n$ $L$  & 
Stores the data in $L$ to a csv file. 
The data is an $m \times n$ matrix in row-major ordering. \\
\hline  
\verb'-mv' & fname-in fname-out  & 
Rename file from fname-in to fname-out. \\
\hline  
\verb'-system' & command  & 
Execute the given system command. \\
\hline  
\verb'-loop' & variable val\_from val\_to val\_step & 
Create a loop with the given loop variable, 
whose values rangle from loop\_from to val\_to minus one in steps of val\_step.  \\
\hline  
\verb'-loop_over' & variable domain & 
Create a loop with the given loop variable, 
whose values rangle over the given domain. \\
\hline  
\verb'-plot_function' & function-name &  \\
\hline  
\verb'-draw_projective_curve' & options &  \\
\hline  
\verb'-tree_draw' & options & 
Draw a tree from a description. 
For options, see Table~\ref{tab:tree:draw}. \\
\hline  
\verb'-extract_from_file' & fname-in key fname-out & 
Extract lines from a text file, 
starting with the given key and continuing until the next empty row.  \\
\hline  
\verb'-extract_from_file_with_tail' & fname-in key tail fname-out  & 
Extract lines from a text file, 
starting with the given key and continuing until the next empty row.  \\
\hline  
\verb'-serialize_file_names' & fname output-mask &  \\
\hline  
\verb'-save_4_bit_data_file' & fname vector-data &  \\
\hline  
\verb'-gnuplot' & fname title label\_x label\_y & 
Produce an diagram of data series using gnuplot~\cite{gnuplot}. \\
\hline  
\verb'-compare_columns' & fname col1 col2 &  \\
\hline  
\verb'-gcd_worksheet' & nb\_problems N key &  \\
\hline  
\verb'-draw_layered_graph' & options &  \\
\hline  
\end{tcolorbox}
\caption{\label{tab:misc:2}Miscellaneous Orbiter Commands (Part 2)}
\end{table}
